\chapter{Fundamentals}

\section{Basic concepts of graph theory}
In this section, we formally introduce the basic concepts of graph theory, mostly following the terminology used in \cite{diestel}.

\begin{definition}
    An \textit{(undirected) graph} is a pair $G = (V, E)$, where $V$ is a non-empty finite set and $E \subseteq {\binom{V}{2}}$ is a set of 2-element subsets of $V$.
    The elements of $V$ are the \textit{vertices} of $G$, and the elements of $E$ are its \textit{edges}. 
    
    A vertex $x$ is \textit{incident} with an edge $e$ if $x \in e$. 
    For an edge $e = \{v, w\}$ we usually just write $vw$ (or $wv$). We say that $e$ \textit{joins} the vertices $v$ and $w$, and $v, w$ are its \textit{ends} or \textit{endvertices}. 
    Furthermore, $v, w$ are called \textit{adjacent} or \textit{neighbours}, and we occasionally denote that by $v \sim w$.

    For any graph $G$ we refer to its vertex set as $V(G)$ and to its edge set as $E(G)$. The number of vertices of a graph $G$ is called its \textit{order}, denoted by $|G|\;:= |V(G)|$. The number of edges is denoted by $e(G) := |E(G)|$.
\end{definition}

Unless explicitly stated, all graphs considered in this thesis are undirected. 
For a directed graph $G$, the definition is essentially the same as for undirected graphs except that $E \subseteq V^2$, i.e. edges are given an ordering and loops are allowed. 
An edge $e = (x, y)$ in $G$ is directed \textit{from} $x$ \textit{to} $y$; $x$ is called its \textit{tail}, and $y$ its \textit{head}.


\begin{example}
    Graphs are typically displayed by drawing a dot for every vertex and connecting adjacent vertices by a line.
    For most purposes, the positioning of the dots and the shape of the lines are not considered important. 
    The relevant information is just which vertices are connected by an edge.
    
    \begin{figure}[H]
        \centering
        \begin{tikzpicture}
            % Define the vertices with manual positions
            \mynode (v1) at (0, 0) [below left] {};
            \mynode (v2) at (1.5, 0.25) {};
            \mynode (v3) at (3, 0) {};
            \mynode (v4) at (-0.5, 1) {};
            \mynode (v5) at (3, 1) {};
            \mynode (v6) at (0.5, 2) {};  % This can be outside for clarity
            \mynode (v7) at (2, 2) {};  % This can be outside for clarity
            
            \node at (v1) [below left] {1};
            \node at (v2) [right  = 0.1cm] {2};
            \node at (v3) [below right] {3};
            \node at (v4) [left  = 0.1cm] {4};
            \node at (v5) [right = 0.1cm] {5};
            \node at (v6) [above  = 0.1cm] {6};
            \node at (v7) [above right] {7};
        
            % Draw the edges
            \draw (v6) -- (v5);
            \draw (v1) -- (v2);
            \draw (v2) -- (v7);
            \draw (v7) -- (v1);
            \draw [-] (v1) .. controls (1.5, -0.75) .. (v3); % Curved line
        \end{tikzpicture}
        \captionsetup{justification=centering, margin=2cm}
        \caption{Graph with vertex set $V = \{1, \dots, 7\}$ and edge set $E = \big\{ \{ 1, 2\}, \{ 1, 3\}, \{ 1, 7\}, \{ 2, 7\}, \{ 5, 6\} \big\}$}
        \label{fig:example-graph}
    \end{figure}
\end{example}

The next example introduces special graphs that are frequently used as examples or counterexamples.
\begin{example} \label{ex:special_graphs}
    \begin{enumerate}
        \item[(a)] The graph $K_n$ on $n$ vertices in which all vertices are pairwise adjacent is called the \textit{complete graph} of order $n$.

        \begin{figure}[H]
            \centering
            \begin{tikzpicture}
                % Define the vertices
                \foreach \i in {1, 2, 3, 4, 5} {
                    \mynode (v\i) at ({360/5 * (\i - 1) + 90} : 1.5) {};
                }
            
                % Draw the edges to form the cycle
                \draw (v1) -- (v2);
                \draw (v1) -- (v3);
                \draw (v1) -- (v4);
                \draw (v1) -- (v5);
                \draw (v2) -- (v3);
                \draw (v2) -- (v4);
                \draw (v2) -- (v5);
                \draw (v3) -- (v4);
                \draw (v3) -- (v5);
                \draw (v4) -- (v5);
                \draw (v5) -- (v1);
            \end{tikzpicture}
    
            \caption{The complete graph $K_5$}
            \label{fig:complete-graph}
        \end{figure}

        \item[(b)] The \textit{path} $P_n$ of length $n$ is a graph of the form 
        \begin{align*}
            V(P_n) = \{x_0, x_1, \dots, x_n\}, && E(P_n) = \{x_0 x_1,\, x_1 x_2,\, \dots,\, x_{n - 1} x_n\} .
        \end{align*}
        The vertices $x_0$ and $x_n$ are called the \textit{ends} of the path, $x_1, \dots, x_{n - 1}$ its \textit{inner vertices}. 

        \item[(c)] The \textit{cycle} $C_n$ of length $n$ is a graph of the form
        \begin{align*}
            V(C_n) = \{x_1, x_2, \dots, x_n\}, && E(C_n) = \{x_1 x_2,\, x_2 x_3,\, \dots,\, x_{n - 1} x_n,\, x_n x_1\} .
        \end{align*}

        \item[(d)] The $n$-dimensional \textit{hypercube} $Q_n$ has the vertex set $V(Q_n) = \{ 0, 1 \}^n$, with two vertices being neighbours if they differ in exactly one coordinate.
        
    \end{enumerate}
\end{example}

Two graphs are defined as equal if and only if they have the same vertex set and the same edge set. 
While this is a completely natural and reasonable definition, it is also quite restrictive: of course, renaming vertices does not really change the internal structure of a graph. 
This motivates the following definition:

\begin{definition}
    Let $G, G'$ be two graphs. An \textit{isomorphism} (from $G$ to $G'$) is a bijection $\varphi: V(G) \rightarrow V(G')$ satisfying 
    \begin{equation*}
        \varphi(x) \varphi(y) \in E(G') \ \Longleftrightarrow  \ xy \in E(G)
    \end{equation*}
    for all $x, y \in V(G)$. If such an isomorphism exists, $G$ and $G'$ are \textit{isomorphic}, and we write $G \cong G'$.
\end{definition}

Graph isomorphism defines an equivalence relation on the class of all graphs. 
Informally speaking, if two graphs are isomorphic, they are equal "up to labelling". 
Most of the time, we are interested in properties of graphs that do not depend on the names of the vertices and do not distinguish between isomorphic graphs.
Therefore, it is justified, for example, to speak of \textit{the} complete graph on $n$ vertices.

\begin{definition}
    Let $G$ be a graph. 
    A graph $H$ is a \textit{subgraph} of $G$ (and $G$ a \textit{supergraph}), denoted $H \subseteq G$, if $V(H) \subseteq V(G)$ and $E(H) \subseteq E(G)$. 
    If $V(H) = V(G)$, $H$ is called \textit{spanning}.
    A subgraph $H$ of $G$ is called \textit{induced} if
    \begin{equation*}
        E(H) = \{ e \in E(G) : e \subseteq V(H) \}.
    \end{equation*}
    An induced subgraph is thus determined by its vertex set, and for any $U \subseteq V(G)$ we denote by $G[U]$ the induced subgraph of $G$ with vertex set $U$. 

    Furthermore, for $U \subseteq V(G)$ we denote by $G - U$ the subgraph induced by $V(G) \setminus U$, i.e. the graph obtained from $G$ by deleting the vertices from $U$ (and all edges incident with $U$). 
    Similarly, if $F$ is a subset of $E(G)$, we denote by $G - F$ the subgraph of $G$ with vertex set $V(G)$ and edge set $E(G) \setminus F$. 
    If $U = \{ v \}$ or $F = \{ e \}$ we simply write $G - v$, respectively $G - e$.
    If $H$ is a subgraph of $G$, we usually write $G \setminus H$ rather than $G - V(H)$.
    Given two non-adjacent vertices $x, y \in V(G)$, $G + xy$ denotes the graph that is obtained from $G$ by adding an edge between $x$ and $y$.
\end{definition}

\begin{definition}
    Let $G$ be a graph.
    A \textit{walk} $W$ (of length $l$) is a sequence $x_0, x_1, \dots, x_l$ of vertices in $G$ such that $x_{i - 1} x_i \in E(G)$ for all $i \in [l]$.
    We refer to such a walk as an $x_0,x_l$-walk.
    If $x_0 = x_l$, the walk is \textit{closed}.

    The walk $W$ is a \textit{path} if all its vertices are pairwise distinct. We say that $x_0, x_l$ are its \textit{endvertices}, and $x_0, x_l$ are \textit{linked} by $W$.
    
    If additionally $l \ge 2$ and $x_l x_0$ is an edge in $G$, then the sequence $x_0, x_1, \dots, x_l, x_0$ is called a \textit{cycle} (of length $l + 1$).
\end{definition}

\begin{definition}
    A graph $G$ is called \textit{bipartite} if there is a partition $V(G) = A \cup B$ of the vertices such that any edge has one endvertex in $A$ and the other in $B$.
\end{definition}

Bipartite graphs are exactly those graphs that do not contain a cycle of odd length (see \cite[Theorem I.4]{bollobas}).

\begin{example}
    The \textit{complete bipartite graph} $K_{m,n}$ is a bipartite graph with partition $V(K_{m, n}) = A \cup B$, $|A|\;=m, |B|\;=n$ where any vertex of $A$ is adjacent to any in $B$.
    That means
    \begin{equation*}
        E(K_{m,n}) = \{ ab : a \in A, b \in B \}.
    \end{equation*}
\end{example}

\begin{definition}
    Let $G$ be a graph. Given vertices $x$ and $y$, their distance $\dist_G (x, y)$ is the length of a shortest $x,y$-path in $G$.
    If no such path exists, then $\dist_G (x, y) := \infty$.

    A graph $G$ is \textit{connected} if any two vertices $x, v \in V(G)$ are linked by an $x,y$-path \mbox{in $G$}.
    A maximal connected subgraph of $G$ is called a \textit{(connected) component}. 
    If a graph $G$ has components $C_1, \dots, C_s$, we also say that $G$ is the disjoint union of its components and write $G = C_1 + \dots + C_s$.
\end{definition}

\begin{definition}
    A connected graph $G$ without any cycles is called a \textit{tree}. 
    A \textit{forest} is a graph whose components are trees.
\end{definition}

\begin{proposition}
        Let $G$ be a graph. The following are equivalent: 
    \begin{enumerate}
        \item $G$ is a tree.
        \item $G$ is a minimal connected graph, i.e., $G$ is connected, but $G - xy$ is disconnected for every edge $xy \in E(G)$. 
        \item $G$ is a maximal acyclic graph, i.e., $G$ does not contain any cycles, but for any two nonadjacent vertices $x, y \in V(G)$, $G + xy$ does contain a cycle.
        \item For any two vertices $x, y \in V(G)$, there is a unique $x, y$-path in $G$.
    \end{enumerate}
\end{proposition}

\begin{proof}
    See \cite[Theorems I.5 and I.6]{bollobas}.
\end{proof}

\begin{corollary}
    Every connected graph contains a spanning tree (i.e. a spanning subgraph that is a tree).
\end{corollary}

% \begin{proof}
%     Every minimal connected spanning subgraph of $G$ is a spanning tree.
% \end{proof}

\begin{definition}
    Let $G$ be a graph. 
    For any vertex $v$, the set of neighbours of $v$ is denoted by $N(v) = N_G(v)$.
    Similarly, for any subset $A \subseteq V(G)$, $N_G(A) := \bigcup_{v \in A} N_G(v)$ denotes the neighbourhood of $A$ in $G$.
        
    The \textit{degree} of a vertex $v$ is defined as \mbox{$d(v) = d_G(v) := |N_G(v)|$}, or equivalently as the number of edges incident with $v$.
    A vertex of degree $0$ is \textit{isolated}.

    The \textit{minimum degree} of a vertex in $G$ is denoted by $\delta(G)$,
    the \textit{maximum degree} \mbox{as $\Delta(G)$}.
    The \textit{average degree} of $G$ is defined as 
    \begin{equation*}
        \bar{d} (G) := \frac{1}{|V(G)|} \sum \limits_{v \in V(G)} d_G (v).
    \end{equation*}
    Clearly, $\delta (G) \le \bar{d} (G) \le \Delta (G)$. 
    A graph in which every vertex has degree $d$ is called $d$-\textit{regular}.    
\end{definition}

The next proposition relates the number of edges with the vertex degrees.
\begin{proposition}
    For every graph $G$, 
    \begin{equation*}
        2 e(G) = \sum_{v \in V(G)} d_G(v).
    \end{equation*}
\end{proposition}

\begin{proof}
    The right-hand side counts every edge twice, namely for both of its endvertices.
\end{proof}

Besides the subgraph relation, the minor relation is another important structural containment relation between two graphs.
\begin{definition}
    Let $G, H$ be graphs. The graph $H$ is a \textit{minor} of $G$ if for each vertex $x \in V(H)$ there exists a subset $V_x \subseteq V(G)$ such that the following holds:
    \begin{enumerate}
        \item $V_x \cap V_y = \emptyset$ for distinct $x, y \in V(H)$.
        \item $G[V_x]$ is connected for any $x \in V(H)$.
        \item For each edge $xy$ in $H$ there exists an edge in $G$ joining $V_x$ and $V_y$.
    \end{enumerate}
    The sets $V_x$ are called \textit{branch sets}.
\end{definition}

The minor relation defines a partial ordering on the class of all (finite) graphs. 
Observe that a subgraph $H$ of a graph $G$ is also its minor:
In this case, the branch sets can simply be chosen as $V_x = \{x\}$ for each $x \in V(H)$.
Figure \ref{fig:example-minor} illustrates the above definition.

\begin{figure}[H]
    \centering
    \begin{tikzpicture}

        % Draw the shapes with a light grey background
        \filldraw[fill=gray!30, draw=black, dashed] (0.25, 0.25) circle [radius=0.6];
        \filldraw[fill=gray!30, draw=black, dashed] (3, 0.25) ellipse [x radius=0.4, y radius=0.6];
        \filldraw[fill=gray!30, draw=black, dashed] (0.2, 2.8) circle [radius=0.6];
        \filldraw[fill=gray!30, draw=black, dashed] (3, 3) circle [radius=0.3];

        
        % Define the vertices with manual positions
        \mynode (v1) at (0, 0) {};
        \mynode (v2) at (0, 0.5) {};
        \mynode (v3) at (0.5, 0.5) {};
        \mynode (v4) at (0.5, 0) {};
        \mynode (v5) at (3, 0) {};
        \mynode (v6) at (3, 0.5) {};
        \mynode (v7) at (3, 3) {};
        \mynode (v8) at (4, 3) {};
        \mynode (v9) at (0, 2.5) {};
        \mynode (v10) at (0, 3) {};
        \mynode (v11) at (0.5, 3) {};

        \mynode (v12) at (-1.5, 1.5) {};
    
        % Draw the edges
        \draw (v1) -- (v2);
        \draw (v2) -- (v3);
        \draw (v3) -- (v4);
        \draw (v4) -- (v1);
        \draw[very thick] (v3) -- (v7);
        \draw[very thick] (v4) -- (v5);
        \draw (v5) -- (v6);
        \draw[very thick] (v6) -- (v7);
        \draw (v7) -- (v8);
        \draw (v9) -- (v10);
        \draw (v10) -- (v11);
        \draw (v11) -- (v9);
        \draw[very thick] (v9) -- (v2);
        \draw[very thick] (v11) -- (v7);
        \draw[very thick] (v9) -- (v6);

        \draw (v12) -- (v9);
        \draw (v12) -- (v1);
        \draw (v12) -- (v10);
    \end{tikzpicture}
    \captionsetup{justification=centering, margin=2cm}
    \caption{A graph with a $K_4$-minor}
    \label{fig:example-minor}
\end{figure}

% An immediate consequence is the so-called Handshaking lemma: 
% \begin{corollary}
%     In every graph, the number of vertices with odd degree is even.
% \end{corollary}

\begin{definition}
    An \textit{independent set} in a graph $G$ is a subset $U \subseteq V(G)$ of the vertices of $G$ none of which are adjacent. 
    The largest size of an independent set in $G$ is called its \textit{independence number} and is denoted by $\alpha(G)$.

    Contrarily, a \textit{clique} in $G$ is a subset $Q \subseteq V(G)$ such that any two vertices in $Q$ are adjacent.
    The largest size of a clique in $G$ is called its \textit{clique number}, denoted $\omega (G)$.
\end{definition}

It is a central activity in graph theory to find sufficient conditions for the occurrence of certain substructures in a graph.
Ramsey's theorem asserts that any graph on sufficiently many vertices contains either a large clique or a large independent set.

\begin{theorem}[Ramsey's theorem] \label{thm:ramsey}
    Let $s, t$ be positive integers. 
    Then there exists a number $R = R(s, t)$ such that any graph $G$ with at least $R$ vertices contains either a clique of size $s$ or an independent set of size $t$.
\end{theorem}

A proof can be found in \cite[Theorem VI.1]{bollobas}, for example.

\section{The matrices associated with a graph}
    \label{section:matrices}
    
We now introduce some of the most prominent matrices studied in spectral graph theory.
Throughout this section, let $G$ be a graph with vertex set $V(G) = \{v_1, \dots, v_n\}$.
The graph $G$ can be represented naturally by its adjacency matrix.

\begin{definition}
    The \textit{adjacency matrix} $A = A_G$ of $G$ is the symmetric $(n \times n)$ matrix defined by
    \begin{equation*}
        A_{ij} = \left\{\begin{array}{ll}
            1, & v_i v_j \in E(G) \\
            0, & \text{otherwise}
        \end{array} \right. .
    \end{equation*}
\end{definition}

\begin{remark}
    In spectral graph theory, it is often convenient to view vectors in $\R^n$ as functions $V(G) \rightarrow \R$, i.e. as vectors in $\R^{V(G)}$.
    Given such a vector $\bm{x}$, we will for the most part simply write $\bm{x}_v$ for the value of $\bm{x}$ at $v$ rather than using integer indices. 
    Similarly, we think about the adjacency matrix $A$ of $G$ as a matrix with rows and columns indexed by $V(G)$, and we write $A_{vw}$ for the entry corresponding to the vertices $v$ and $w$.
    The adjacency matrix then defines a linear map $\R^{V(G)} \rightarrow \R^{V(G)}$.
    When we multiply a vector $\bm{x}$ to the adjacency matrix, we obtain
    \begin{equation*}
        (A \bm{x})_v = \sum \limits_{w \sim v} \bm{x}_w,
    \end{equation*}
    which means that the resulting value at a vertex $v$ is the sum of the values of $\bm{x}$ at the neighbours \mbox{of $v$}.
\end{remark}

The next proposition states that the powers of the adjacency matrix encode the number of walks of a given length.

\begin{proposition} \label{prop:count-walks}
    Let $m$ be a nonnegative integer. 
    The entry $(A^m)_{vw}$ is the number of walks of length $m$ from $v$ to $w$. 
    In particular, $(A^2)_{vv}$ is the degree of the vertex $v$, and $\tr A^2 = 2 e(G)$.
\end{proposition}

\begin{proof}
    We use induction on $m$. For $m = 0$ we have $A^0 = I$, and the statement is true. (Note that a walk of length 0 consists of a single vertex.)
    Now suppose the claim holds for a fixed $m \in \N_0$.
    Then
    \begin{equation*}
        (A^{m + 1})_{vw} = \sum \limits_{u \in V(G)} (A^m)_{vu} \cdot A_{uw}.
    \end{equation*}
    By the induction hypothesis, $(A^m)_{vu}$ equals the number of $v, u$-walks of length $m$. 
    Furthermore, $A_{uw} = 1$ if and only if there is an edge from $u$ to $w$.
    Hence $(A^{m + 1})_{vw}$ counts the number of $v, u$-walks of length $m$ joined by an edge $u w$, which provides exactly the number of $v, w$-walks of length $m + 1$.
\end{proof}

% \begin{definition}
%     The \textit{incidence matrix} $B = B_G$ of $G$ is the $(n \times m)$-matrix defined by
%     \begin{equation*}
%         B_{ij} = 
%         \begin{cases}
%             1,\; v_i \in e_j \\
%             0,\; \text{otherwise}
%         \end{cases}.
%     \end{equation*}
% \end{definition}

% By the definition of the matrix multiplication we have the following connection between the adjacency matrix and the incidence matrix of $G$.
% \begin{proposition}
%     $B B^T = A + D$, 
%     where $D$ is the diagonal matrix with $D_{ii} = d_G(v_i)$.
% \end{proposition}

Another important matrix that arises in the context of spectral graph theory is the Laplace matrix of $G$.
\begin{definition}
    The \textit{Laplace matrix} or \textit{Laplacian} $L = L_G$ of $G$ is defined by $L = D - A$, where $D$ is the degree diagonal matrix of $G$, i.e. the diagonal matrix with $D_{vv} = d(v)$. 
    Hence, the entries of the Laplace matrix are given by
    \begin{equation*}
        L_{vw} = \left\{\begin{array}{ll}
            d(v),& v = w \\
            -1, & vw \in E(G) \\
            0,& \text{otherwise}
        \end{array} \right. .
    \end{equation*}
\end{definition}

\begin{definition}
    An \textit{orientation} of $G$ is the assignment of a direction to every edge, i.e., any undirected edge $e = \{x, y \}$ is replaced by a directed edge $(x, y)$ (or $(y, x)$). 

    Fix an orientation of $G$. The \textit{directed incidence matrix} $N = N_G$ of $G$ with respect to the given orientation is the ($|V(G)| \times |E(G)|$) matrix defined by
    \begin{equation*}
        N_{ve} = \left\{\begin{array}{ll}
            1,& v \text{ is the head of } e \\
            -1, & v \text{ is the tail of } e \\
            0,& \text{otherwise}
        \end{array} \right. ,
    \end{equation*}
    where the rows are indexed by the vertices and the columns are indexed by the edges \mbox{of $G$}.
\end{definition}

\begin{proposition} \label{prop:laplacian-quadratic-form}
    Let $N$ denote the directed incidence matrix of $G$ w.r.t. to some orientation. 
    Then $L = N N^T$, and for any vector $\bm{v} \in \R^{V(G)}$,
    \begin{equation*}
        \bm{v}^T L \bm{v} = \sum \limits_{x \sim y} (\bm{v}_x - \bm{v}_y)^2,
    \end{equation*}
    where the sum is over all edges $xy \in E(G)$. 
    In particular, $L$ is positive semidefinite.
\end{proposition}

\begin{proof}
    By the definition of the matrix multiplication, 
    \begin{equation*}
        (N N^T)_{xy} = \sum \limits_{e \in E(G)} N_{xe} N_{ye}.
    \end{equation*}
    For the summands we have 
    \begin{equation*}
        N_{xe} N_{ye} = \left\{\begin{array}{ll}
            1, & x = y \text{ and } x \in e \\
            -1, & e = \{ x, y \} \\
            0, & \text{otherwise }
        \end{array} \right. .
    \end{equation*}
    Hence, 
    \begin{equation*}
         (N N^T)_{xy} = \left\{\begin{array}{ll}
            d_G(x), & x = y \\
            -1, & x \sim y \\
            0, & \text{otherwise}
        \end{array} \right \} = L_{xy}.
    \end{equation*}
    Now let $\bm{v}$ be a vector. Then for any edge $e = \{x, y\}$, $(N^T \bm{v})_e = \pm (\bm{v}_x - \bm{v}_y)$.
    It follows that 
    \begin{equation*}
        \bm{v}^T L \bm{v} = (\bm{v}^T N) (N^T \bm{v}) = \sum \limits_{e \in E(G)} (N^T \bm{v})_e = \sum \limits_{x \sim y} (\bm{v}_x - \bm{v}_y)^2.
    \end{equation*}
\end{proof}

There is an abundance of literature about the Laplace matrix of a graph and its eigenvalues (for an introduction, see e.g. \cite{brouwerhaemers} or \cite[Chapter 13]{godsilroyle}).
In this thesis, we will focus on the properties of the adjacency matrix of a graph, though the Laplace matrix will play a part in the proof of Cheeger's inequality (\mbox{Theorem \ref{thm:cheeger}}).

\section{The spectrum of a graph}
Let $G$ be an undirected graph with $n$ vertices. 
Since the adjacency matrix $A_G$ is symmetric, the spectral theorem of linear algebra states that all its $n$ eigenvalues are real and there exists an orthonormal basis of $\R^{V(G)}$ consisting of eigenvectors of $A_G$. 
Furthermore, the algebraic and geometric multiplicities of each eigenvalue coincide, so we will simply speak about \textit{multiplicity}. 

\begin{definition}
    The \textit{spectrum} of $G$ is defined as the spectrum of $A_G$, that is, the set of all eigenvalues of $A_G$ together with their multiplicities.
    The \textit{spectral radius} of $G$ is the maximum absolute value of an eigenvalue of $A_G$.
\end{definition}

\begin{remark}
    While the adjacency matrix $A_G$ is determined by the ordering of the vertices, the spectrum of $G$ is independent of the representation chosen. 
    Indeed, any two adjacency matrices $A$, $A'$ of $G$ are similar because $A' = P^T AP$ for some permutation matrix $P \in \{0, 1\}^{n \times n}$ (and $P^T = P^{-1}$).
\end{remark}

For any eigenvalue $\lambda$ of $A_G$ with a corresponding eigenvector $\bm{x}$, the eigenvalue equation $A_G \bm{x} = \lambda \bm{x}$ can be equivalently stated as 
\begin{equation*}
    \lambda \bm{x}_v = \sum \limits_{u \sim v} \bm{x}_u, \ \ v \in V(G).
\end{equation*}
In accordance with the definition of the spectrum of $G$, we refer to the eigenvalues and eigenvectors of $A_G$ as the eigenvalues and eigenvectors of $G$.
Similarly, the characteristic polynomial of $G$, denoted $p_G$, is that of $A_G$. 
% Since $p_G$ is a polynomial with integer coefficients, the eigenvalues of a graph are totally real algebraic integers. 
We will denote the eigenvalues of $G$ by $\lambda_1(G), \dots, \lambda_n(G)$, or simply $\lambda_1, \dots, \lambda_n$ if there is no possibility of confusion. 
By convention, the eigenvalues are always assumed to be in a decreasing order:
\begin{equation*}
    \lambda_1 \ge \lambda_2 \ge \dots \ge \lambda_n.
\end{equation*}
Recall that the trace of a matrix is equal to the sum of its eigenvalues (counted according to their algebraic multiplicity).
In a graph $G$ without self-loops, the diagonal entries of $A_G$ are all zero, and therefore 
\begin{equation*}
    \sum \limits_{i = 1}^n \lambda_i(G) = \tr A_G = 0.
\end{equation*}

As already mentioned, $G$ has eigenvectors $\bm{x}_1, \dots, \bm{x}_n$ (such that $A_G \bm{x}_i = \lambda_i \bm{x}_i$) that form an orthonormal basis of $\R^{V(G)}$.
This gives rise to the \textit{spectral decomposition} 
\begin{equation*}
A_G = \sum _{i = 1}^n \lambda_i \bm{x}_i \bm{x}_i ^T
\end{equation*}
of $A_G$. 
To see why this equation holds, let $B = \sum _{i = 1}^n \lambda_i \bm{x}_i \bm{x}_i ^T$. 
Since $\bm{x}_1, \dots, \bm{x}_n$ is a basis of $\R^{V(G)}$, it suffices to show that $(A_G - B) \bm{x}_j = 0$ for all $j \in [n]$.
Indeed, we have
\begin{equation*}
    (A_G - B) \bm{x}_j = A_G \bm{x}_j - \sum _{i = 1}^n \lambda_i \bm{x}_i \bm{x}_i ^T \bm{x}_j = \lambda_j \bm{x}_j - \lambda_j \bm{x}_j = 0.
\end{equation*}
For the second equality, we have used the fact that the $\bm{x}_i$ are orthonormal.

Interestingly, many structural properties of a graph are encoded in its spectrum. 
Before exploring these connections, we will first consider the spectra of some special graphs. 
% On this occasion, we quote the spectral mapping theorem for polynomials.
%
% \begin{theorem}[Spectral mapping theorem]
%     Let $A \in \R^{n \times n}$ be a symmetric matrix, and let $\sigma(A)$ denote its spectrum.
%     Then, for any polynomial $p(X) \in \R[X]$,
%     \begin{equation*}
%         \sigma(p(A)) = \{ p(\lambda) : \lambda \in \sigma (A) \}.
%     \end{equation*}
% \end{theorem}
% A proof can be found in \cite[Theorem 7.36]{nairsinghlinearalgebra}, for instance.
%
We will denote eigenvalue multiplicities greater than one as superscripts, i.e., the notation $\lambda^{(k)}$ means that the eigenvalue $\lambda$ occurs with multiplicity $k$.

\begin{example} \label{ex:special_spectra}
    \begin{enumerate}
        \item[(a)] The empty graph $G = ([n], \emptyset)$ on $n$ vertices has \mbox{spectrum $0^{(n)}$}.
        \item[(b)] The spectrum of the complete graph $K_n$ is $(n - 1), -1^{(n - 1)}$. 
            Indeed, $K_n$ has the adjacency matrix $A = J - I$, and $J$ has spectrum $n, 0^{(n - 1)}$.
        \item[(c)] The spectrum of the complete bipartite graph $K_{m,n}$ is $\sqrt{mn}, 0^{(m + n - 2)}, -\sqrt{mn}$.
            To see that, consider the adjacency matrix of $K_{m,n}$, which is of the form
            \begin{equation*}
                A = \begin{bmatrix}
                    0 & J_{m,n} \\ 
                    J_{n, m} & 0
                \end{bmatrix}.
            \end{equation*}
            As $A$ has rank 2, it has eigenvalue $0$ with multiplicity $m + n - 2$. 
            Using the eigenvectors
            \begin{equation*}
                (\underbrace{\sqrt{n}, \dots, \sqrt{n}}_{m \text{ times}}, \underbrace{\pm \sqrt{m}, \dots, \pm \sqrt{m}}_{n \text{ times}})^T
            \end{equation*}
            one can verify that $\pm \sqrt{mn}$ are eigenvalues.

        \item[(d)] The spectrum of the cycle $C_n$ consists of the numbers $2 \cos (\frac{2 \pi k}{n})$, $k = 0, \dots, n-1$.
            The adjacency matrix of $C_n$ has the form
            \begin{equation*}
                A = \begin{bmatrix}
                    0 & 1 & \dots & 0 & 1 \\
                    1 & 0 & 1 & & 0 \\ 
                    \vdots & 1 & 0 & \ddots & \vdots \\ 
                    0 & &\ddots & \ddots & 1 \\
                    1 & 0 & \dots & 1 & 0 
                \end{bmatrix}.
            \end{equation*}
            For any $n$-th root of unity $\xi \in \C$, $(1, \xi, \xi^2, \dots, \xi^{n - 1})^T$ is an eigenvector of $A$ for the eigenvalue $\xi + \xi^{-1}$.
            The $n$-th roots of unity are $e^{2 \pi i k /n}$, $k = 0, \dots, n - 1$, and since $e^{2 \pi i k / n} + e^{-2 \pi i k/n} = 2 \cos (2 \pi k / n)$, we obtain the claimed spectrum.

        \item[(e)] The $n$-dimensional hypercube $Q_n$ has the eigenvalues $n - 2k$ with multiplicity $\binom{n}{k}$, $k = 0, \dots, n$. 
            We show this inductively. For $n = 1$, we have $Q_1 \cong K_2$, hence the eigenvalues are $\pm 1$ by (b).
            Now, for a fixed $n$, let $A_n$ denote the adjacency matrix of $Q_n$. 
            It is not hard to see that $Q_{n + 1}$ has an adjacency matrix of the form 
            \begin{equation*}
                A_{n + 1} = \begin{bmatrix}
                    A_n & I \\ 
                    I & A_n 
                \end{bmatrix},
            \end{equation*}
            where the $A_n$ blocks correspond to $n$-dimensional subcubes \mbox{of $Q_{n + 1}$}.
            Let $\bm{v}$ be an eigenvector of $A_n$ for an eigenvalue $\lambda$. Then
            \begin{equation*}
                \begin{bmatrix}
                    A_n & I \\ 
                    I & A_n 
                \end{bmatrix}
                \begin{bmatrix}
                    \bm{v} \\ \bm{v} 
                \end{bmatrix}
                = (\lambda + 1)
                \begin{bmatrix}
                    \bm{v} \\\bm{v} 
                \end{bmatrix}
            \end{equation*}
            and 
            \begin{equation*}
                \begin{bmatrix}
                    A_n & I \\ 
                    I & A_n 
                \end{bmatrix}
                \begin{bmatrix}
                    \bm{v} \\ -\bm{v} 
                \end{bmatrix}
                = (\lambda - 1)
                \begin{bmatrix}
                    \bm{v} \\ -\bm{v} 
                \end{bmatrix}.
            \end{equation*}
            Thus, an eigenvalue $\lambda$ of $Q_n$ yields eigenvalues $\lambda \pm 1$ of $Q_{n + 1}$.
            By the induction hypothesis it follows that $Q_{n + 1}$ has eigenvalues 
            \begin{equation*}
                n + 1, n - 1, n - 3, \dots, -n - 1.
            \end{equation*}
            Since the eigenvectors given above are all linearly independent, the multiplicity of $n + 1 - 2k$ ($k = 0, \dots, n + 1$) is obtained by adding the multiplicities of $n - 2(k - 1)$ and $n - 2k$ as eigenvalues in $Q_n$:
            \begin{equation*}
                \binom{n}{k - 1} + \binom{n}{k} = \binom{n + 1}{k},
            \end{equation*}
            where $\binom{n}{-1} = \binom{n}{n + 1} := 0$. As $\sum_{k = 0}^{n + 1} \binom{n + 1}{k} = 2^{n + 1}$, these are indeed all eigenvalues.
    \end{enumerate}
\end{example}

The next proposition tells us that the spectrum of a disconnected graph consists of the spectra of its components.
\begin{proposition}
    Let $G$ be a graph with connected components $C_1, \dots, C_s$. 
    Then the spectrum of $G$ is the union of the spectra of the $C_i$, with multiplicities being added.
\end{proposition}

\begin{proof}
    $G$ admits an adjacency matrix of the form
    \begin{equation*}
        A = \begin{bmatrix}
            A_{C_1} & 0 & \dots & 0 \\
            0 & A_{C_2} & & \vdots \\
            \vdots & & \ddots & 0 \\
            0 & \dots & 0 & A_{C_s}
        \end{bmatrix}.
    \end{equation*}
    Hence, the characteristic polynomial of $G$ satisfies $p_G (t) = p_{C_1} (t) \cdot \dots \cdot p_{C_s} (t)$, and the claim follows.
\end{proof}

It is worth noting that while the spectrum is a graph invariant (i.e., it is the same for isomorphic graphs), graphs are not generally determined by their spectrum.
Indeed, there exist pairs of graphs that are not isomorphic but share the same spectrum. Such graphs are called \textit{cospectral}. 
For example, we can infer from Example \ref{ex:special_spectra} that $K_{1, 4}$ and $K_1 + C_4$ both have the spectrum $2, 0^{(3)}, -2$.

\section{Rayleigh quotients}
Consider a symmetric matrix $A \in \R^{n \times n}$. For any vector $\bm{x} \in \R^n \setminus \{0\}$, the term 
\begin{equation*}
    \frac{\bm{x}^T A \bm{x}}{\bm{x}^T \bm{x}}
\end{equation*}
is called the \textit{Rayleigh quotient} of $\bm{x}$ with respect to $A$.
The next theorem tells us that the eigenvalues of a symmetric matrix can be characterised as the maxima (resp. minima) of the Rayleigh quotients over certain subspaces. 

\begin{theorem}[Rayleigh] \label{thm:rayleigh}
    Let $A \in \R^{n \times n}$ be a symmetric matrix with eigenvalues $\lambda_1 \ge \dots \ge \lambda_n$.
    Let $\bm{v}_1, \dots, \bm{v}_n$ be orthonormal eigenvectors of $A$ (such that $A \bm{v}_i = \lambda_i \bm{v}_i$).
    For $k \in [n]$, let $S_k = \langle \bm{v}_1, \dots, \bm{v}_k \rangle$. 
    Then
    \begin{equation*}
        \lambda_k = \min \limits_{\bm{x} \in S_k \setminus \{0 \}} \frac{\bm{x}^T A \bm{x}}{\bm{x}^T \bm{x}} = \min \limits_{\bm{x} \in S_k, \|\bm{x}\| = 1} \bm{x}^T A \bm{x}
    \end{equation*}
    and
    \begin{equation*}
        \lambda_k = \max \limits_{\bm{x} \in S_{k - 1}^{\perp} \setminus \{0 \}} \frac{\bm{x}^T A \bm{x}}{\bm{x}^T \bm{x}} = \max \limits_{\bm{x} \in S_{k - 1}^{\perp}, \| \bm{x} \| = 1} \bm{x}^T A \bm{x}.
    \end{equation*}
    Moreover, a vector $\bm{x}$ attains the maximum (resp. minimum) if and only if $\bm{x}$ is a $\lambda_k$-eigenvector \mbox{of $A$}.
\end{theorem}

This result is sometimes also called the Courant--Fischer theorem.
We reserve this name for a related, but slightly different result (see e.g. \cite[Theorem 4.2.6]{matrixanalysis}) and refer to the above as Rayleigh's principle.

% \begin{proof}
%     Let $\bm{x} \in \R^n, \bm x \ne 0$ and write
%     \begin{equation*}
%         \bm x = \sum_{i = 1}^n \alpha_i \bm{v}_i
%     \end{equation*}
%     with suitable coefficients $\alpha_i \in \R$. Since the $\bm{v}_i$ are orthonormal eigenvectors, 
%     \begin{equation*}
%         \bm{x}^T \bm{x} = \sum_{i = 1}^n \alpha_i^2 \ \ \text{ and } \ \ \bm{x}^T A \bm{x} = \sum_{i = 1}^n \alpha_i^2 \lambda_i.
%     \end{equation*}
%     If $\bm{x} \in S_k$, then $\alpha_{k + 1} = \dots = \alpha_n = 0$, hence
%     \begin{equation*}
%         \frac{\bm{x}^T A \bm{x}}{\bm{x}^T \bm{x}} = \frac{\alpha_1^2 \lambda_1 + \dots + \alpha_k ^2 \lambda_k}{\alpha_1^2 + \dots + \alpha_k ^2}
%             \ge \frac{(\alpha_1^2 + \dots + \alpha_k ^2) \lambda_k}{\alpha_1^2 + \dots + \alpha_k ^2} = \lambda_k.
%     \end{equation*}
%     Similarly, if $\bm{x} \in S_{k - 1}^{\perp}$, then $\alpha_{1} = \dots = \alpha_{k - 1} = 0$, hence
%     \begin{equation*}
%         \frac{\bm{x}^T A \bm{x}}{\bm{x}^T \bm{x}} = \frac{\alpha_k^2 \lambda_k + \dots + \alpha_n ^2 \lambda_n}{\alpha_k^2 + \dots + \alpha_n ^2}
%             \le \frac{(\alpha_k^2 + \dots + \alpha_n ^2) \lambda_k}{\alpha_k^2 + \dots + \alpha_n ^2} = \lambda_k.
%     \end{equation*}
%     In both cases, equality holds if and only $\bm{x}$ is a $\lambda_k$-eigenvector. 
%     If $\bm{x} \ne 0$, then $\bm{\xi} = \bm{x} / \|\bm{x}\|$ is a unit vector with 
%     \begin{equation*}
%         \bm{\xi}^T A \bm{\xi} = \frac{\bm{x}^T A \bm{x}}{\bm{x}^T \bm{x}},
%     \end{equation*}
%     so it does not matter whether the optimization is over all nonzero vectors or over unit vectors only.
% \end{proof}

\begin{proof}
    We only show the first characterisation, as the other one is completely analogous.
    Let $\bm{x} \in \R^n \setminus \{ 0 \}$, and expand $\bm{x}$ in the orthonormal basis of eigenvectors of $A$, that is,
    \begin{equation*}
        \bm x = \sum_{i = 1}^n \alpha_i \bm{v}_i
    \end{equation*}
    with suitable coefficients $\alpha_i \in \R$. Since the $\bm{v}_i$ are orthonormal eigenvectors, 
    \begin{equation*}
        \bm{x}^T \bm{x} = \sum_{i = 1}^n \alpha_i^2 \ \ \text{ and } \ \ \bm{x}^T A \bm{x} = \sum_{i = 1}^n \alpha_i^2 \lambda_i.
    \end{equation*}
    Assume $\bm{x} \in S_k$. 
    Then $\alpha_{k + 1} = \dots = \alpha_n = 0$, and hence
    \begin{equation*}
        \frac{\bm{x}^T A \bm{x}}{\bm{x}^T \bm{x}} = \frac{\alpha_1^2 \lambda_1 + \dots + \alpha_k ^2 \lambda_k}{\alpha_1^2 + \dots + \alpha_k ^2}
            \ge \frac{(\alpha_1^2 + \dots + \alpha_k ^2) \lambda_k}{\alpha_1^2 + \dots + \alpha_k ^2} = \lambda_k.
    \end{equation*}
    If $\bm{x}$ is a $\lambda_k$-eigenvector, then the above holds with equality.
    
    Conversely, if the vector $\bm{x}$ satisfies $\bm{x}^T A \bm{x} / \bm{x}^T \bm{x} = \lambda_k$, we must have $\alpha_i = 0$ for all $i \le k$ with $\lambda_i > \lambda_k$.
    This implies that $\bm{x}$ is the linear combination of $\lambda_k$-eigenvectors, and therefore a $\lambda_k$-eigenvector itself.
\end{proof}

The most useful characterisations are summarised in the following corollary.
\begin{corollary} \label{cor:rayleigh-quotients}
    Let $A \in \R^{n \times n}$ be a symmetric matrix with eigenvalues $\lambda_1 \ge \dots \ge \lambda_n$ and let $\bm{v}_1$ be a $\lambda_1$-eigenvector. Then
    \begin{align*}
        \lambda_1 = \max \limits_{\bm{x} \ne 0} \frac{\bm{x}^T A \bm{x}}{\bm{x}^T \bm{x}}, &&
        \lambda_2 = \max \limits_{\bm{x} \ne 0, \bm{x} \perp \bm{v}_1} \frac{\bm{x}^T A \bm{x}}{\bm{x}^T \bm{x}}, &&
        \lambda_n = \min \limits_{\bm{x} \ne 0} \frac{\bm{x}^T A \bm{x}}{\bm{x}^T \bm{x}}.
    \end{align*}
\end{corollary}



% \begin{proof}
%     Since there is only one zero-dimensional subspace of $\R^n$, namely $S = \{ 0 \}$, we can omit the outer minimum in (\ref{eq:cf2}) for $i = 0$ and obtain the formula for $\lambda_1$.
%     Similarly, we can just take $S = \R^n$ in (\ref{eq:cf1}) for $i = n$ and obtain the formula for $\lambda_n$.
    
%     For $\lambda_2$, we have
%     \begin{equation*}
%         \lambda_2 \le \max \limits_{0 \ne x, x \perp v_1} \frac{x^T A x}{x^T x}
%     \end{equation*}
%     by (\ref{eq:cf2}). 
%     For the other direction, let $x_0$ be the maximizer of $\frac{x^T A x}{x^T x}$ among all $x \ne 0, x \perp v_1$. 
%     Then we can write $x_0 = \sum _{i = 2}^n \alpha_i v_i$ for suitable $\alpha_i \in \R$. Using the orthonormality of the $v_i$ it follows that 
%     \begin{align*}
%         x_0^T A x_0 &= \bigg( \sum \limits_{i, j = 2}^n \alpha_i v_i \bigg) A \bigg( \sum \limits_{i = 2}^n \alpha_i v_i \bigg) 
%             = \bigg( \sum \limits_{i = 2}^n \alpha_i v_i \bigg) \bigg(\sum \limits_{i = 2}^n \alpha_i \lambda_i v_i \bigg) \\
%         &= \sum_{i = 2}^n \lambda_i \alpha_i^2 \le \lambda_2 \sum_{i = 2}^n  \alpha_i^2 = \lambda_2 \|x_0\|,
%     \end{align*}
%     and we conclude
%     \begin{equation*}
%         \max \limits_{x \ne 0, x \perp v_1} \frac{x^T A x}{x^T x} = \frac{x_0^T A x_0}{x_0^T x_0} \le \frac{\lambda_2 \|x_0\|}{\|x_0\|} = \lambda_2.
%     \end{equation*}
%     This proves the formula for $\lambda_2$.
% \end{proof}

Using Rayleigh's principle, we can obtain basic information about the magnitude of the largest eigenvalue $\lambda_1$ of a graph.
\begin{proposition}
    For every graph $G$, $\bar{d} (G) \le \lambda_1 \le \Delta(G)$.
    If $G$ is connected, then $\lambda_1 = \Delta(G)$ if and only if $G$ is $\Delta(G)$-regular.
\end{proposition}

\begin{proof}
    We follow the proof given in \cite{spielman}.
    As usual, let $A$ denote the adjacency matrix of $G$.
    To show the lower bound, consider the Rayleigh quotient of the all-ones vector $\bm{1}$:
    \begin{equation*}
        \lambda_1 = \max \limits_{\bm{x} \ne 0} \frac{\bm{x}^T A \bm{x}}{\bm{x}^T \bm{x}} 
            \ge \frac{\bm{1}^T A \bm{1}}{\bm{1}^T \bm{1}}
            = \frac{\sum_{u, v} A_{uv}}{n} 
            = \frac{2 e(G)}{n} = \bar{d} (G).
    \end{equation*}
    For the upper bound, let $\bm \psi$ be a $\lambda_1$-eigenvector and let $v \in V(G)$ be a vertex at which $\bm \psi (v)$ is maximal.
    Without loss of generality, we can assume that $\bm \psi (v) > 0$, otherwise \mbox{consider $-\bm{\psi}$}. 
    Then 
    \begin{equation*}
        \lambda_1 = \frac{(A \bm{\psi})_v}{\bm{\psi}_v} 
            = \frac{\sum_{u \sim v} \bm{\psi}_u}{\bm{\psi}_v} 
            = \sum \limits_{u \sim v} \frac{\bm{\psi}_u}{\bm{\psi}_v} 
            \le \sum \limits_{u \sim v} 1 = d(v) \le \Delta(G). 
            \tag{$\star$}
    \end{equation*}

    For the second part of the proposition, assume that $G$ is connected. 
    If $G$ is $\Delta(G)$-regular, then $A \bm{1} = \Delta(G) \bm{1}$, so $\Delta(G) = \lambda_1$.
    For the other direction, assume that that ($\star$) holds with equality.
    It follows that \mbox{$d(v) = \Delta(G)$}, and $\bm{\psi}_u = \bm{\psi}_v$ for every vertex $u$ adjacent to $v$.
    Hence, in the equation ($\star$) we may replace the vertex $v$ by any of its neighbours $u$ and apply the same reasoning to $u$.
    By iterating the argument and using that $G$ is connected, we conclude that $d(u) = \Delta(G)$ for all $u \in V(G)$.
    Thus, $G$ is $\Delta(G)$-regular.
\end{proof}

Rayleigh's principle together with the spectral decomposition of the adjacency matrix of a graph can also be used to deduce lower and upper bounds on the number of closed walks.
\begin{lemma} \label{lem:bounding-walks}
    Let $G$ be a graph, and let $\bm{\psi}$ denote a normalised $\lambda_1$-eigenvector of $G$.
    For each vertex $v$ and $k \in \N$, the number of closed walks of length $2k$ rooted at $v$, namely $(A^{2k})_{vv}$, can be estimated as follows:
    \begin{equation*}
        \bm{\psi}_v^2 \lambda_1 ^{2k} \le (A^{2k})_{vv} \le \lambda_1 ^{2k}.
    \end{equation*}
\end{lemma}

\begin{proof}
    For the first inequality, consider an orthonormal set $\bm{x}_1 = \bm{\psi}, \bm{x}_2, \dots, \bm{x}_n$ of eigenvectors of $G$, and recall the spectral decomposition of $A$: 
    \begin{equation*}
        A = \sum \limits_{i = 1}^n \lambda_i \bm{x}_i \bm{x}_i^T.
    \end{equation*}
    Since $A^{2k}$ has the eigenvalues $\lambda_i^{2k}$ with the same eigenvectors as $A$, the spectral decomposition for $A^{2k}$ is  
    \begin{equation*}
        A^{2k} = \sum \limits_{i = 1}^n \lambda_i^{2k} \bm{x}_i \bm{x}_i^T.
    \end{equation*}
    Hence,
    \begin{equation*}
        (A^{2k})_{vv} = \sum \limits_{i = 1}^n \lambda_i^{2k} \underbrace{(\bm{x}_i \bm{x}_i^T)_{vv}}_{= \bm{x}_{i, v}^2} \ge \bm{\psi}_v^2 \lambda_1^{2k}.
    \end{equation*}
    The second inequality follows directly from Rayleigh's principle: 
    \begin{equation*}
        (A^{2k})_{vv} = \bm{e}^T_v A^{2k} \bm{e}_v \le \lambda_1^{2k}.
    \end{equation*}
\end{proof}

\section{Interlacing}

Let $0 < m < n$ be integers, and consider two sequences $\lambda_1 \ge \dots \ge \lambda_n$ and $\mu_1 \ge \dots \ge \mu_m$ of real numbers. 
Following \cite{brouwerhaemers}, we say that the second sequence \textit{interlaces} the first one if 
\begin{equation*}
    \lambda_i \ge \mu_i \ge \lambda_{n - m + i}, \ \ i \in [m].
\end{equation*}
If $m = n - 1$, these inequalities become
\begin{equation*}
    \lambda_1 \ge \mu_1 \ge \lambda_2 \ge \mu_2 \ge \dots \ge \mu_{m} \ge \lambda_{n}.
\end{equation*}
The interlacing is \textit{tight} if there exists $k \in \{0, 1, \dots, m \}$ such that 
\begin{equation*}
    \lambda_i = \mu_i, \ \ i \in \{ 1, \dots, k\} \ \text{ and } \ \mu_i = \lambda_{n -m + i}, \ \ i \in \{ k + 1, \dots, m\}.
\end{equation*}

The following is a classical result and a consequence of Rayleigh's principle.
\begin{proposition}[Cauchy interlacing theorem]
    Let $A \in \R^{n \times n}$ be a symmetric matrix, and let $B \in \R^{n \times n}$ be a principal submatrix of $A$. 
    Then the eigenvalues of $B$ interlace those \mbox{of $A$}.
\end{proposition}

\begin{proof}
    Without loss of generality, we can assume that $B$ consists of the first $m$ rows and columns \mbox{of $A$}. 
    (Otherwise permutate the indices of $A$ accordingly, which leaves the spectrum invariant.)
    Let $\lambda_1 \ge \dots \ge \lambda_n$ denote the eigenvalues of $A$ and let $\bm{v}_1, \dots, \bm{v}_n$ be a corresponding set of orthonormal eigenvectors.
    Furthermore, let $\mu_1 \ge \dots \ge \mu_m$ denote the eigenvalues of $B$, with respective eigenvectors $\bm{w}_1, \dots, \bm{w}_m$.
    We have $B = S^T A S$ and $S^T S = I_m$, where 
    \begin{equation*}
        S = \begin{bmatrix}
            I_m \\ 0
        \end{bmatrix} \in \R^{n \times m}.
    \end{equation*}
    Fix $i \in [m]$, and choose a nonzero vector 
    \begin{equation*}
        \bm{x} \in \langle \bm{w}_1, \dots, \bm{w}_i \rangle \cap \Big \langle S^T \bm{v}_1, \dots, S^T \bm{v}_{i - 1} \Big \rangle ^\perp.
    \end{equation*}
    Such a vector exists because %$\dim \left \langle S^T \bm{v}_1, \dots, S^T \bm{v}_{i - 1} \right \rangle \le i - 1$, and hence
    \begin{equation*}
        \dim \langle \bm{w}_1, \dots, \bm{w}_i \rangle + \dim \Big \langle S^T \bm{v}_1, \dots, S^T \bm{v}_{i - 1} \Big \rangle ^\perp \ge i + (m - i + 1) > m.
    \end{equation*}
    Then $S \bm x \in \langle \bm{v}_1 \dots, \bm{v}_{i - 1} \rangle^\perp$, and by Rayleigh's principle, it follows that 
    \begin{equation*}
        \lambda_i \ge \frac{(S \bm{x})^T A (S\bm{x})}{(S \bm{x})^T (S \bm{x})} 
            = \frac{\bm{x}^T (S^T A S) \bm{x}}{\bm{x}^T (S^T S) \bm{x}}
            = \frac{\bm{x}^T B \bm{x}}{\bm{x}^T \bm{x}} \ge \mu_i.
    \end{equation*}
    % For the other inequality, consider 
    % \begin{equation*}
    %     \bm{x} \in \langle \bm{w}_i, \dots, \bm{w}_m \rangle \cap \Big \langle S^T \bm{v}_{n - m + i + 1}, \dots, S^T \bm{v}_{n} \Big \rangle ^\perp.
    % \end{equation*}
    % Then $\bm x \in \langle \bm{w}_1, \dots, \bm{w}_{i - 1} \rangle ^\perp$ and $$S \bm x \in \langle \bm{v}_1 \dots, \bm{v}_{i - 1} \rangle^\perp$$
    For the other inequality, the same argument carried out on $-A$ and $-B$ yields 
    \begin{equation*}
        -\lambda_{n + 1 - k} \ge -\mu_{m + 1 - k},\ \ k \in [m].
    \end{equation*}
    Substituting $k = m + 1 - i$, we conclude that $\mu_i \ge \lambda_{n - m + i}$.
\end{proof}

In graph-theoretic terminology, the Cauchy interlacing theorem translates to:
\begin{corollary}
    Let $H$ be an induced subgraph of a graph $G$. Then the eigenvalues of $H$ interlace those of $G$.
\end{corollary}

\begin{proof}
    The adjacency matrix $A_H$ of $H$ is a principal submatrix of $A_G$.
\end{proof}

The following bound on the smallest eigenvalue of a graph is a simple application of interlacing. 
\begin{proposition}
    Let $G$ be a graph with at least one edge and smallest \mbox{eigenvalue $\lambda_n$}.
    If $G$ is the disjoint union of complete graphs, then $\lambda_n = -1$. Otherwise $\lambda_n \le -\sqrt{2}$.
\end{proposition}

\begin{proof}
    Recall that the smallest eigenvalue of a complete graph with at least one edge is $-1$, and the spectrum of a disconnected graph is the union of the spectra of its components.
    Hence, if $G$ is the disjoint union of complete graphs, its smallest eigenvalue \mbox{is $-1$}. 
    Otherwise, there are (distinct) vertices $x, y, z$ such that $xy, yz \in E(G)$ and $xz \notin E(G)$, i.e., $G$ has an induced $K_{1,2}$.
    The smallest eigenvalue of $K_{1, 2}$ is $-\sqrt{2}$, and by the Cauchy interlacing theorem we conclude that $-\sqrt{2} \ge \lambda_n$.
\end{proof}

Besides principal submatrices, interlacing can also be applied to so-called quotient matrices. 
Let $A \in \R ^{n \times n}$ be a symmetric matrix whose rows and columns are partitioned according to a partition $\{ X_1, \dots, X_m \}$ of $[n]$, that is, 
\begin{equation*}
    A = \begin{bmatrix}
        A_{1,1} & \dots & A_{1,m} \\
        \vdots & & \vdots \\
        A_{m,1} & \dots & A_{m,m} \\
    \end{bmatrix},
\end{equation*}
where $A_{i,j}$ is the submatrix of $A$ consisting of the rows in $X_i$ and the columns \mbox{in $X_j$}.
Let $$b_{i,j} := \frac{1}{|X_i|} \bm{1}^T A_{i,j}\bm{1}$$ denote the average row sum \mbox{of $A_{i,j}$}.
The \textit{quotient matrix} of $A$ with respect to the given partition is the matrix
\begin{equation*}
    B = \begin{bmatrix}
        b_{1,1} & \dots & b_{1,m} \\
        \vdots & & \vdots \\
        b_{m,1} & \dots & b_{m,m} \\
    \end{bmatrix} \in \R^{m \times m}.
\end{equation*}
If the row sums within each block $A_{i,j}$ are constant, that is $A_{i,j} \bm{1} = b_{i,j} \bm{1}$ for all $i,j \in [m]$, the partition is called \textit{equitable}. 
In the situation described, we have the following result.

\begin{proposition} \label{prop:partition-interlacing}
    The eigenvalues of $B$ interlace those of $A$. If the interlacing is tight, the partition is equitable.
\end{proposition}
For a proof, see \cite[Corollary 2.5.4]{brouwerhaemers}.

\section{The Perron--Frobenius theorem} \label{section:perron-frobenius}

The Perron--Frobenius theorem is one of the most important tools in spectral graph theory. 
We will give a version of the result as in \cite{brouwerhaemers}, though it is stated in more detail there.

For matrices $A, B \in \R^{n \times n}$, the notation $A \le B$ means that 
\begin{equation*}
    A_{ij} \le B_{ij}, \ \ i, j \in [n].
\end{equation*}
If all of the above inequalities are strict, we write $A < B$. 
In particular, $A \ge 0$ (or $A > 0$) means that all entries of $A$ are nonnegative (resp. positive).
We shall use the same notation for vectors.

\begin{definition}
    Let $A \in \R^{n \times n}$, $A \ge 0$.
    The matrix $A$ is called \textit{irreducible} if for all $i,j \in [n]$ there exists an integer $k$ such that $(A^k)_{ij} > 0$.
\end{definition}

\begin{theorem}[Perron--Frobenius] \label{thm:perron-frobenius}
    Let $A \in \R^{n \times n}$, $A \ge 0$ be an irreducible matrix. 
    Then there is a unique number $\lambda_0 > 0$ with the following properties:
    \begin{enumerate}
        \item $\lambda_0$ is an eigenvalue of $A$ with algebraic multiplicity 1.
        \item For each eigenvalue $\lambda$ of $A$, we have $|\lambda|\, \le \lambda_0$. 
        \item There is a vector $\bm{v} > 0$ such that $A \bm{v} = \lambda_0 \bm{v}$.
            Besides $\lambda_0$, no other eigenvalue of $A$ has a nonnegative eigenvector.  
            % Furthermore, if $A$ has precisely $m$ eigenvalues $\lambda$ with $|\lambda|\, = \lambda_0$, then these are obtained by multiplying $\lambda_0$ with the $m$-th roots of unity, namely
            % $\lambda = \lambda_0 \cdot e^{\frac{2 \pi i k}{m}}$ for $k = 0, \dots, m - 1$.
        % \item No other eigenvalue of $A$ besides $\lambda_0$ has a nonnegative eigenvalue.
        % \item Let $t > 0$, and let $\bm{x} \in \R^n \setminus \{0 \}$, $\bm{x} \ge 0$ be a vector. 
        %     If $A \bm{x} \le t \bm{x}$, then $t \ge \lambda_0$. 
        %     If $A \bm{x} \ge t \bm{x}$, then $t \le \lambda_0$.
        %     In both cases, equality holds if and only if $A \bm{x} = t \bm{x}$.
        \item Let $B \in \R^{n \times n}$ with $0 \le B \le A$ be another matrix. 
            If $\mu$ is an eigenvalue of $B$, then $|\mu|\, \le \lambda_0$. 
            Furthermore, if $|\mu|\, = \lambda_0$, then $B = A$.
            The same holds if $B$ is a principal submatrix of $A$.
    \end{enumerate}
\end{theorem}
A proof can be found in \cite[Theorem 2.2.1]{brouwerhaemers}.
% \begin{proof}
%     See \cite[Theorem 2.2.1]{brouwerhaemers}.
% \end{proof}
For more details on Perron--Frobenius theory, see also \cite[Chapter XIII]{gantmacher}.

The Perron--Frobenius theorem immediately provides information about the largest eigenvalue of a graph and the corresponding eigenvector.
\begin{proposition} \label{prop:perron-frobenius}
    Let $G$ be a graph with largest eigenvalue $\lambda_1$.
    \begin{enumerate}
        \item We have $|\lambda|\; \le \lambda_1$ for each eigenvalue $\lambda$ of $G$.
    \end{enumerate}
    Furthermore, if $G$ is connected, then:
    \begin{enumerate}[resume]
        \item $\lambda_1$ has multiplicity $1$, and a corresponding eigenvector can be chosen to have only positive entries.
        \item If $H$ is a proper subgraph of $G$, then $\lambda_1 (H) < \lambda_1$.
    \end{enumerate}
\end{proposition}

\begin{proof}
    Observe that $A_G \ge 0$, and by Proposition \ref{prop:count-walks} $G$ being connected is equivalent to $A_G$ being irreducible.
    Hence, the Perron--Frobenius theorem immediately implies (ii) and (iii).
    To show (i), consider any component $C$ of $G$.
    By the Perron--Frobenius theorem, we have $|\lambda|\, \le \lambda_1 (C)$ for each eigenvalue $\lambda$ of $C$.
    Since the spectrum of $G$ is the union of the spectra of its components, (i) follows.
\end{proof}

Due to the various implications of the Perron--Frobenius theorem, the largest eigenvector $\lambda_1$ of a graph is also referred to as the \textit{Perron eigenvalue} of $G$, and a corresponding eigenvector is called a \textit{Perron eigenvector}.
Observe that if $G$ is disconnected, a Perron eigenvector $\bm x$ can still be chosen to have only nonnegative entries by letting $\bm x$ be strictly positive on a component with spectral radius $\lambda_1$ and zero on all other components.

The next application of the Perron--Frobenius theorem gives a characterisation of connected bipartite graphs.
\begin{proposition}
    Let $G$ be a connected graph. Then the following are equivalent:
    \begin{enumerate}
        \item $G$ is bipartite.
        \item The spectrum of $G$ is symmetric about $0$, i.e., for each eigenvalue $\lambda$ of $G$, $-\lambda$ is an eigenvalue as well, with the same multiplicity.
        \item $-\lambda_1(G)$ is an eigenvalue of $G$.
    \end{enumerate}
\end{proposition}

We adopt the proof from \cite[Theorem 8.8.2]{godsilroyle}.
\begin{proof}
    We start by showing (i) $\Rightarrow$ (ii).
    Clearly, if $G$ is bipartite, it has an adjacency matrix of the form 
    \begin{equation*}
        A = \begin{bmatrix} 
                  0 & B \\ 
                  B^T & 0
                \end{bmatrix}
    \end{equation*}
    for a suitable $\{ 0, 1\}$-matrix $B$. Now if the partitioned vector $(\bm{v}, \bm{w})$ is an eigenvector for the eigenvalue $\lambda$, $(\bm{v}, -\bm{w})$ is an eigenvector for the eigenvalue $-\lambda$, and (ii) follows.

    The implication (ii) $\Rightarrow$ (iii) is trivial, so it remains to show (iii) $\Rightarrow$ (i).
    Assume that $-\lambda_1$ is an eigenvalue of $G$, and let $A$ denote the adjacency matrix of $G$.
    Consider the graph $G^2$ induced by $A^2$, i.e. with vertex set $V(G)$ and with an edge between two vertices $x$ and $y$ whenever $(A^2)_{xy} \ne 0$ and $x \ne y$.
    Since $A^2$ has the eigenvalue $\lambda_1 ^2$ with multiplicity at least 2, $G^2$ is disconnected by the Perron--Frobenius theorem.
    Observe that a walk of length $2l$ in $G$ corresponds to a walk of length $l$ in $G^2$, 
    so two vertices are in the same connected component of $G^2$ if and only if they are linked by a walk of even length in $G$.
    Since $G$ is connected and $G^2$ is not, there are vertices $x, y$ such that every $x, y$-walk in $G$ has odd length. 
    Let $P$ be such a walk. 
    We claim that $G$ does not contain an odd cycle:
    Indeed, otherwise, we would find an $x, y$-walk of even length in $G$ by walking from $x$ to a vertex of that odd cycle, walking around the cycle once, returning to $x$ on the same way and walking along $P$ to $y$. 
    We conclude that $G$ is bipartite.
\end{proof}
